\chapter{Mechanistic Models of Preparative Chromatography}
\label{ch:models}

\begin{quote}\itshape
``All models are wrong, but some are useful.'' --- G.~E.~P.~Box, whose name also
graces the experimental designs of Chapter~\ref{ch:design}.
\end{quote}

\noindent
A chromatographic column separates the components of a mixture by exploiting the
differences in how strongly each component partitions between a flowing
\emph{mobile phase} and a stationary, porous \emph{adsorbent}. Feed is loaded,
washed, and then eluted---often by changing the composition of the mobile phase
so as to weaken binding---and the components emerge at the outlet at different
times. To \emph{model} this process is to write down, and then solve, the
conservation laws that govern the migration of each component through the bed,
closed by a constitutive law (the \emph{adsorption isotherm}) that says how much
of each component is bound at equilibrium as a function of the local mobile-phase
state.

This chapter builds that model from the ground up. Section~\ref{sec:governing}
derives the governing equations. Section~\ref{sec:engine} specialises them to the
form the package integrates and describes the numerical method.
Section~\ref{sec:isotherm} introduces the unifying idea that organises the whole
package---\emph{one isotherm core, and the chromatography mode is nothing but the
law by which a modulator controls binding affinity}. Sections~\ref{sec:iex}
through~\ref{sec:rp} treat the individual modes. Section~\ref{sec:program}
describes how an elution program (in particular, a gradient) is specified, and
Section~\ref{sec:metrics} defines the performance metrics that turn a simulated
chromatogram into the numbers a process developer cares about.
Section~\ref{sec:legacy} presents the reduced equilibrium-dispersive model the
package retains for fast single-component work, and Section~\ref{sec:grm} the
\emph{general rate model}---a second, finite-volume solver that resolves the
film and intraparticle-diffusion mass transfer explicitly and is mass
conservative where the lumped engine is not. We close, in
Section~\ref{sec:limits-summary}, with a consolidated statement of the model
family's limits of validity.

%==============================================================================
\section{The column as a distributed-parameter system}
\label{sec:governing}

Consider a cylindrical column of length $L$ and internal diameter $d$, packed
with porous adsorbent particles. Let the \emph{bed porosity} (the volume
fraction occupied by the flowing interstitial liquid) be $\eps$. A
mobile-phase pump drives liquid through the bed at an \emph{interstitial
velocity} $u$ (the true velocity of the liquid between particles, related to the
volumetric flow rate $Q$ and the empty cross-sectional area $A=\pi d^2/4$ by
$Q = u A \eps$). These quantities, and the derived volumes, are exactly the
fields of the package's \code{ColumnGeometry} data structure:
%
\begin{equation}
A = \frac{\pi d^2}{4}, \qquad
V_{\mathrm{col}} = A L, \qquad
V_{\mathrm{resin}} = V_{\mathrm{col}}(1-\eps).
\label{eq:geom}
\end{equation}

We track, for each component $i$, two concentrations defined per unit volume:
$c_i(z,t)$, the concentration in the mobile phase, and $q_i(z,t)$, the
concentration adsorbed on the stationary phase (expressed per unit volume of
adsorbent). A mass balance on a thin slice of the column, $[z, z+\dd z]$,
equates the rate of accumulation to the net flux in by convection and
dispersion, minus the rate at which material is captured by the solid:
%
\begin{equation}
\underbrace{\eps\,\pdv{c_i}{t}}_{\text{mobile accumulation}}
+ \underbrace{(1-\eps)\,\pdv{q_i}{t}}_{\text{uptake by solid}}
= \underbrace{-\,\eps u \pdv{c_i}{z}}_{\text{convection}}
+ \underbrace{\eps D_{\mathrm{ap}}\,\pdv[2]{c_i}{z}}_{\text{axial dispersion}} .
\label{eq:fullbalance}
\end{equation}
%
Here $D_{\mathrm{ap}}$ is an \emph{apparent axial dispersion coefficient}
($\mathrm{m^2\,s^{-1}}$): a lumped term that represents, collectively, molecular
diffusion, the spread of residence times caused by the tortuous flow path
(eddy dispersion), and---crucially---the band broadening that arises because mass
transfer into and out of the particles is not instantaneous. Dividing through by
$\eps$ and introducing the \emph{phase ratio}
%
\begin{equation}
\phi \;=\; \frac{1-\eps}{\eps},
\label{eq:phaseratio}
\end{equation}
%
the volume of solid per unit volume of liquid, gives the convection--dispersion
equation in the form the package uses:
%
\begin{equation}
\pdv{c_i}{t}
= -\,u\,\pdv{c_i}{z}
+ D_{\mathrm{ap}}\,\pdv[2]{c_i}{z}
- \phi\,\pdv{q_i}{t}.
\label{eq:cdmobile}
\end{equation}

Equation~\eqref{eq:cdmobile} is one equation in two unknowns, $c_i$ and $q_i$.
It must be closed by a statement of how the adsorbed phase responds to the
mobile phase. There are two classical closures, and the package's design hinges
on the choice between them.

\paragraph{The equilibrium closure.}
If mass transfer is fast compared with convection, the two phases are everywhere
in local equilibrium, $q_i = q_i^{\mathrm{eq}}(c)$, and one may substitute
$\partial q_i/\partial t = (\dd q_i/\dd c_i)\,\partial c_i/\partial t$ to fold the
isotherm into a single \emph{retardation factor}. This is the classical
\emph{equilibrium-dispersive} (ED) model; it is fast but cannot, by
construction, represent a mobile-phase composition that changes in time without
introducing awkward source terms. The package keeps an ED model for
single-component, isocratic work (Section~\ref{sec:legacy}).

\paragraph{The rate closure.}
If, instead, one writes a \emph{kinetic} law for the approach of $q_i$ to its
equilibrium value, the adsorbed phase becomes an independent state variable that
can be carried explicitly. This is the \emph{transport-dispersive} model, and it
is the general engine of the package. It is the rate closure that makes
\emph{gradient elution}---the central technique of preparative
chromatography---natural to simulate, because the mobile-phase modulator may
vary in time without any $\dd H/\dd t$ source term appearing.

%==============================================================================
\section{The transport-dispersive engine}
\label{sec:engine}

The package's general model (\code{models/chromatography/engine.py}) augments the
state at each point of the column with the adsorbed concentrations $q_i$ and with
a \emph{modulator} $m(z,t)$---the mobile-phase variable that governs binding
strength (salt concentration for ion exchange and hydrophobic interaction; the
volume fraction of organic solvent for reversed phase). The modulator is treated
as an unretained tracer: it is convected and dispersed but does not adsorb.

The full coupled system, for $n_c$ components, is
%
\begin{align}
\text{modulator:}\quad
&\pdv{m}{t} = -\,u\,\pdv{m}{z} + D_{\mathrm{ap}}\,\pdv[2]{m}{z},
\label{eq:mod}\\[0.4em]
\text{mobile:}\quad
&\pdv{c_i}{t} = -\,u\,\pdv{c_i}{z} + D_{\mathrm{ap}}\,\pdv[2]{c_i}{z}
  - \phi\,\pdv{q_i}{t},
\label{eq:mobile}\\[0.4em]
\text{stationary:}\quad
&\pdv{q_i}{t} = k_{m,i}\,\bigl(\qstar_i(\bm c, m, \mathrm{pH}) - q_i\bigr).
\label{eq:ldf}
\end{align}
%
Equation~\eqref{eq:ldf} is the \emph{linear driving force} (LDF) approximation:
the rate of adsorption is proportional to the distance of the adsorbed
concentration from its equilibrium value $\qstar_i$, with a lumped
\emph{mass-transfer coefficient} $k_{m,i}$ (units $\mathrm{s^{-1}}$). The
equilibrium loading $\qstar_i$ is supplied by the isotherm of
Section~\ref{sec:isotherm}.

\begin{keybox}{Why carry the adsorbed phase explicitly?}
Folding the isotherm into a retardation factor (the ED model) requires the
mobile-phase composition to be constant, because the retardation factor then
depends on time and a $\partial H/\partial t$ source term appears. By keeping
$q_i$ as a state variable governed by the rate law~\eqref{eq:ldf}, the inlet
modulator $m(0,t)$ may follow an arbitrary gradient with \emph{no} extra terms.
The price is $n_c$ additional fields. The mass-transfer coefficient $k_{m,i}$
then does double duty: it is a physical rate constant, and it is the knob that
sets the band-broadening (and hence the plate count) of the simulated peaks. The
equilibrium-dispersive limit is recovered as $k_{m,i}\to\infty$.
\end{keybox}

\subsection{Boundary and initial conditions}

The column is fed at the inlet ($z=0$) and discharges at the outlet ($z=L$):
%
\begin{align}
\text{inlet (Dirichlet):}\quad & c_i(0,t) = c_{i,\mathrm{feed}}(t), \qquad
  m(0,t) = m_{\mathrm{in}}(t), \\
\text{outlet (Neumann):}\quad & \left.\pdv{c_i}{z}\right|_{z=L} = 0, \qquad
  \left.\pdv{m}{z}\right|_{z=L} = 0.
\end{align}
%
The inlet feed concentration $c_{i,\mathrm{feed}}(t)$ is a rectangular pulse
during the load phase and zero otherwise; the inlet modulator $m_{\mathrm{in}}(t)$
follows the elution program (Section~\ref{sec:program}). The zero-gradient
outlet condition is the physically appropriate statement that the column
discharges into open tubing. Initially the column is pre-equilibrated at the
starting modulator value and is free of protein:
$c_i(z,0)=0$, $q_i(z,0)=0$, $m(z,0)=m_{\mathrm{in}}(0)$.

\subsection{Spatial discretisation: the method of lines}
\label{sec:mol}

The column is divided into $N$ finite-volume cells of width $\Delta z = L/N$
(the field \code{n\_cells}). The spatial derivatives are replaced by finite
differences, converting the partial differential
equations~\eqref{eq:mod}--\eqref{eq:ldf} into a large system of ordinary
differential equations in time---one per state variable per cell. The schemes
are chosen for stability rather than formal order of accuracy:

\begin{center}
\renewcommand{\arraystretch}{1.35}
\begin{tabular}{@{}p{0.27\linewidth} p{0.36\linewidth} p{0.29\linewidth}@{}}
\toprule
\textbf{Term} & \textbf{Scheme} & \textbf{Rationale} \\
\midrule
Convection $u\,\partial c/\partial z$ &
First-order \emph{upwind}: $u(c_i - c_{i-1})/\Delta z$ &
Stable and non-oscillatory at every P\'eclet number \\
Dispersion $D_{\mathrm{ap}}\,\partial^2 c/\partial z^2$ &
Second-order \emph{central}: $D_{\mathrm{ap}}(c_{i+1}-2c_i+c_{i-1})/\Delta z^2$ &
Accurate and stable for $D_{\mathrm{ap}}>0$ \\
Inlet boundary & Dirichlet ghost cell $c_0 = c_{\mathrm{feed}}(t)$ &
Direct imposition of the feed \\
Outlet boundary & Zero-gradient ghost cell $c_{N+1}=c_N$ &
Column discharges to open tubing \\
\bottomrule
\end{tabular}
\end{center}

\paragraph{Numerical dispersion.}
First-order upwinding introduces an artificial diffusion of magnitude
$D_{\mathrm{num}} \approx u\,\Delta z/2$. In practice this is folded into the
interpretation of $D_{\mathrm{ap}}$: the simulated band width reflects the sum of
the physical dispersion and the grid-induced dispersion, so the cell count $N$
becomes, with $k_{m,i}$, part of how the plate count is tuned. A converged
solution requires $\Delta z$ small enough that $D_{\mathrm{num}}$ is comfortably
below the physical $D_{\mathrm{ap}}$.

\subsection{Time integration: implicit BDF}

The resulting ODE system is \emph{stiff}. The stiffness has two sources: the
rate constant $k_{m,i}$ can be large (fast local equilibration), and, when
binding is strong, a small change in mobile concentration drives a large change
in adsorbed concentration through the $-\phi\,\partial q_i/\partial t$ coupling
in~\eqref{eq:mobile}. An explicit integrator would be forced to take
ruinously small time steps. The package therefore uses the implicit
\emph{Backward Differentiation Formula} (BDF) integrator of
\code{scipy.integrate.solve\_ivp}, which selects its order and step size
adaptively from the local dynamics and is unconditionally stable for this class
of problem.

\paragraph{Exploiting sparsity.}
An implicit integrator must repeatedly form (or approximate) the Jacobian of the
right-hand side. Formed by dense finite differences this would cost $\mathcal
O(n)$ right-hand-side evaluations per step, where $n = 2 n_c N + N$ is the total
state dimension. The package instead hands the solver the \emph{sparsity
pattern} of the Jacobian (\code{\_jacobian\_sparsity}). The structure is local:
in cell $j$, the mobile concentration $c_{i,j}$ couples only to its transport
neighbours $c_{i,j\pm1}$, to every component's concentration in the
\emph{same} cell (through the competitive isotherm), to its own $q_{i,j}$, and to
the modulator $m_j$; the modulator couples only to its transport neighbours.
Supplying this pattern lets the solver reconstruct the Jacobian from a handful of
graph-coloured evaluations instead of $n$, which is the dominant speed-up that
makes a gradient run feasible inside a design loop.

\begin{remark}
The modulator is floored at a small positive value ($10^{-9}$) before it is used
in the isotherm. This guards the ion-exchange affinity law, which contains the
factor $(\Lambda/m)^{\nu}$ and would otherwise be singular should dispersion
overshoot drive the local salt concentration momentarily negative.
\end{remark}

%==============================================================================
\section{One isotherm core; the mode is the modulator law}
\label{sec:isotherm}

The constitutive heart of the model is the equilibrium isotherm $\qstar_i(\bm c,
m, \mathrm{pH})$, supplied by \code{models/chromatography/isotherms.py}. The
package's organising insight is that \emph{every} common chromatography mode can
be written with the \emph{same} multi-component competitive-Langmuir functional
form, and that the modes differ \emph{only} in how a single scalar
affinity factor depends on the modulator. Concretely,
%
\begin{equation}
\boxed{\;
\qstar_i = \frac{q_{\max,i}\,b_i(m,\mathrm{pH})\,c_i}
                {1 + \sum_{j=1}^{n_c} b_j(m,\mathrm{pH})\,c_j}
\;}
\qquad\text{(nonlinear / overload),}
\label{eq:langmuir}
\end{equation}
%
where $q_{\max,i}$ is the saturation capacity of component $i$ (g per litre of
adsorbent) and $b_i$ is its dimensionless affinity. In the \emph{dilute} limit,
where the bound amount is far below saturation
($\sum_j b_j c_j \ll 1$), the denominator collapses to unity and the isotherm
becomes linear,
%
\begin{equation}
\qstar_i = H_i(m,\mathrm{pH})\, c_i, \qquad
H_i = q_{\max,i}\, b_i(m,\mathrm{pH}),
\label{eq:henry}
\end{equation}
%
with $H_i$ the \emph{Henry constant} (the initial slope of the isotherm). The
\code{Isotherm} object carries a Boolean \code{linear} flag that selects between
the two forms; the competition denominator in~\eqref{eq:langmuir} is precisely
what allows components to displace one another on an overloaded column, and is
therefore what makes multi-component resolution possible.

\begin{keybox}{The unifying picture}
Choose a mode $\Leftrightarrow$ choose the affinity law $b_i(m,\mathrm{pH})$.
Everything else---transport, numerics, metrics---is shared. The three laws below
(Figure~\ref{fig:affinity}) are the entire mode-specific content of the model.
\end{keybox}

\begin{figure}[t]
\centering
\includegraphics[width=\linewidth]{affinity_laws.png}
\caption[The three modulator affinity laws]{The affinity factor $b(m)$ for the
three modes, on a logarithmic ordinate. \emph{Left:} ion exchange (SMA)---binding
falls steeply as salt rises, $b\propto m^{-\nu}$. \emph{Centre:} hydrophobic
interaction (salting-out)---binding \emph{grows} exponentially with salt.
\emph{Right:} reversed phase (linear solvent strength)---binding falls
exponentially as the organic fraction rises. The sign and shape of these curves
dictate the direction of the elution gradient in each mode.}
\label{fig:affinity}
\end{figure}

%==============================================================================
\section{Ion-exchange chromatography (CEX / AEX)}
\label{sec:iex}

\subsection{The Steric-Mass-Action affinity law}

In ion-exchange chromatography the adsorbent bears fixed charges that are
initially neutralised by small counter-ions (typically from the salt). A charged
protein binds by displacing several of these counter-ions; elution is achieved by
raising the salt concentration, which shifts the competition back toward the
counter-ion. The \emph{Steric-Mass-Action} (SMA) formalism captures this
stoichiometric exchange. In the package the SMA law enters through the affinity
factor (class \code{SMALaw}):
%
\begin{equation}
b_i(m,\mathrm{pH})
= \beta_i
\left(\frac{\Lambda}{m}\right)^{\nu_i}
\exp\!\bigl[\nu_{\mathrm{pH},i}\,(\mathrm{pH}-\mathrm{pH}_{\mathrm{ref}})\bigr],
\label{eq:sma}
\end{equation}
%
so that, in the dilute limit~\eqref{eq:henry}, the Henry constant is
%
\begin{equation}
H_i = q_{\max,i}\,\beta_i
\left(\frac{\Lambda}{m}\right)^{\nu_i}
\exp\!\bigl[\nu_{\mathrm{pH},i}\,(\mathrm{pH}-\mathrm{pH}_{\mathrm{ref}})\bigr].
\label{eq:smahenry}
\end{equation}
%
The two dependences this law encodes are the two great levers of an
ion-exchange step. \emph{Salt} enters as $m^{-\nu_i}$: because the
characteristic charge $\nu_i$ is typically $3$--$8$ for a protein, a modest
increase in salt produces an order-of-magnitude collapse in binding---this is
what makes a salt gradient such a powerful elution tool. \emph{pH} enters
through the exponential factor: it tunes the protein's net charge and hence its
affinity. A \emph{cation} exchanger (\code{cation\_exchange}) takes
$\nu_{\mathrm{pH}}<0$, so binding strengthens as the pH falls below the reference
and the protein becomes more positively charged; an \emph{anion} exchanger
(\code{anion\_exchange}) takes $\nu_{\mathrm{pH}}>0$. The two builders are
literally the same function with opposite sign of $\nu_{\mathrm{pH}}$.

\subsection{Parameters}

\begin{center}
\renewcommand{\arraystretch}{1.3}
\begin{longtable}{@{}l l p{0.52\linewidth}@{}}
\toprule
\textbf{Symbol} & \textbf{Units} & \textbf{Meaning} \\
\midrule
\endhead
$\beta_i$ & --- & Affinity prefactor of component $i$; sets the absolute binding
   strength at the reference conditions (\code{beta}). \\
$\nu_i$ & --- & Characteristic charge: the number of counter-ions displaced on
   binding; governs salt sensitivity (\code{nu}). \\
$\Lambda$ & mM & Ionic capacity of the resin: the concentration of fixed charges
   available for exchange (\code{ionic\_capacity}). \\
$\nu_{\mathrm{pH},i}$ & --- & Sensitivity of $\ln b_i$ to pH; negative for CEX,
   positive for AEX (\code{nu\_ph}). \\
$\mathrm{pH}_{\mathrm{ref}}$ & --- & Reference pH at which the pH factor is unity
   (\code{ph\_ref}). \\
$q_{\max,i}$ & g\,L$^{-1}$ & Saturation capacity, per litre of adsorbent; the
   overload ceiling (\code{q\_max}). \\
$m$ & mM & The modulator: here the mobile-phase salt concentration. \\
\bottomrule
\end{longtable}
\end{center}

Figure~\ref{fig:henry} shows the resulting Henry-constant surface
$H(m,\mathrm{pH})$ for a representative cation exchanger: binding spans several
decades across the operating window, and the salt and pH axes are the two
critical process parameters that an ion-exchange characterisation study sets out
to map.

\begin{figure}[t]
\centering
\includegraphics[width=0.72\linewidth]{henry_surface.png}
\caption[SMA Henry-constant surface]{The base-10 logarithm of the SMA Henry
constant, $\log_{10} H(m,\mathrm{pH})$, for a model cation exchanger
($\nu=4$, $\nu_{\mathrm{pH}}=-1.2$, $\Lambda=1000$~mM). Binding falls steeply
toward high salt (the elution direction) and weakens as pH rises (the protein
loses positive charge). The near-linear, widely spaced contours are the
signature of the power-law/exponential form of Eq.~\eqref{eq:smahenry}.}
\label{fig:henry}
\end{figure}

\subsection{The elution chromatogram}

Loaded at low salt, where $H$ is enormous and the protein is effectively
immobile, and then eluted by a rising salt gradient, the protein migrates only
once the local salt has climbed enough to bring $H$ toward unity. The result is
the characteristic gradient-elution peak of Figure~\ref{fig:cex}: a sharp band
emerging at a well-defined salt concentration, the \emph{elution salt}, which is
itself a reproducible, parameter-dependent fingerprint of the molecule.

\begin{figure}[t]
\centering
\includegraphics[width=0.85\linewidth]{cex_gradient.png}
\caption[CEX gradient chromatogram]{A simulated cation-exchange linear
salt-gradient separation (solid: outlet protein concentration; dashed: the inlet
salt gradient), computed by \code{run\_column}. The protein is retained through
load and wash at low salt and elutes as a compact band once the gradient lifts
the salt past its critical elution value.}
\label{fig:cex}
\end{figure}

\begin{appbox}{ion-exchange chromatography}
Ion exchange is the most widely used preparative mode in the manufacture of
therapeutic proteins. \emph{Cation} exchange is the standard \emph{polishing}
step after Protein~A capture of a monoclonal antibody, where it clears
aggregates, host-cell proteins, leached Protein~A, and charge variants, and is
frequently run in \emph{bind-and-elute} mode with a salt or pH gradient.
\emph{Anion} exchange is the classic \emph{flow-through} polishing step for
antibodies (the product flows past while acidic impurities, DNA, and endotoxin
bind) and a capture step for many acidic proteins, plasmids, and viral vectors.
The salt/pH dependence encoded by Eq.~\eqref{eq:sma} is exactly the design space
a chromatography development scientist explores when setting the load pH and the
gradient.
\end{appbox}

\begin{databox}{ion exchange (SMA)}
The SMA parameters are obtained from a small, classical set of experiments.
\textbf{(i) Linear-gradient elution (LGE) experiments.} A series of dilute
injections is run at several gradient slopes; the salt concentration at peak
apex, $m_R$, is recorded. Plotting $\log(\text{normalised gradient slope})$
against $\log m_R$ yields a straight line---the \emph{Yamamoto plot}---whose
slope and intercept give the characteristic charge $\nu_i$ and the equilibrium
prefactor $\beta_i$ directly. \textbf{(ii) Isocratic retention.} Measuring the
retention factor at several constant salt concentrations gives $\nu_i$ from the
slope of $\log k'$ versus $\log m$. \textbf{(iii) pH scouting.} Repeating the
above at several pH values calibrates $\nu_{\mathrm{pH},i}$ and
$\mathrm{pH}_{\mathrm{ref}}$. \textbf{(iv) Breakthrough curves and batch
isotherms.} Frontal analysis (loading to saturation and recording the
breakthrough) yields the saturation capacity $q_{\max,i}$ and the resin ionic
capacity $\Lambda$ (the latter also by acid--base or counter-ion titration).
Inverse fitting of these data to the model is exactly the task of the package's
uncertainty-quantification layer.
\end{databox}

\begin{limitbox}{ion exchange (SMA)}
The implemented SMA law is the \emph{dilute-limit} affinity~\eqref{eq:sma}
embedded in a competitive-Langmuir denominator. Relative to the full implicit
SMA isotherm, in which the available charge is reduced by the steric shielding
$\sigma_i q_i$ of already-bound molecules, this form omits the explicit steric
term: the overload ceiling is represented by $q_{\max,i}$ rather than by the
$(\Lambda - \sum(\nu_i+\sigma_i)q_i)$ factor of the rigorous theory. It is
therefore most accurate when the column is not driven deep into overload. The law
also assumes a single, fixed ionic capacity $\Lambda$ and a pH effect that is
log-linear over the window of interest; it does not represent conformational
changes, charge regulation, or salt-mediated changes in $q_{\max}$. Finally, the
characteristic charge $\nu_i$ is treated as constant, whereas for some proteins
it drifts with pH.
\end{limitbox}

%==============================================================================
\section{Hydrophobic-interaction chromatography (HIC)}
\label{sec:hic}

Hydrophobic-interaction chromatography exploits the opposite salt dependence.
The adsorbent presents mildly hydrophobic ligands; at \emph{high} salt the
ordered water around exposed hydrophobic patches on the protein is disrupted
(the ``salting-out'' effect), and the protein binds. Elution is achieved by
\emph{lowering} the salt. The package encodes this with the salting-out affinity
law (class \code{SaltingOutLaw}):
%
\begin{equation}
b_i(m) = \beta_i \, \exp\!\bigl(K_{s,i}\, m\bigr),
\label{eq:hic}
\end{equation}
%
with $\beta_i$ the affinity at zero salt and $K_{s,i}$ the \emph{salting-out
coefficient} (units mM$^{-1}$). Because $b_i$ grows exponentially with salt
(Figure~\ref{fig:affinity}, centre), a HIC step is loaded at high salt and eluted
by a \emph{descending} salt gradient---the mirror image of ion exchange.

\begin{center}
\renewcommand{\arraystretch}{1.3}
\begin{tabular}{@{}l l p{0.55\linewidth}@{}}
\toprule
\textbf{Symbol} & \textbf{Units} & \textbf{Meaning} \\
\midrule
$\beta_i$ & --- & Residual affinity extrapolated to zero salt (\code{beta}). \\
$K_{s,i}$ & mM$^{-1}$ & Salting-out coefficient; the slope of $\ln b_i$ with
   salt, set by the protein's surface hydrophobicity and the salt's position in
   the Hofmeister series (\code{ks}). \\
$q_{\max,i}$ & g\,L$^{-1}$ & Saturation capacity (\code{q\_max}). \\
$m$ & mM & Modulator: mobile-phase salt (e.g.\ ammonium sulfate). \\
\bottomrule
\end{tabular}
\end{center}

\begin{appbox}{hydrophobic-interaction chromatography}
HIC is a powerful \emph{polishing} and \emph{aggregate-clearance} step. Because
it separates on hydrophobicity rather than charge, it is orthogonal to ion
exchange and is often paired with it in a platform sequence to attack impurities
that co-elute on a charge basis---most importantly antibody aggregates, which
expose more hydrophobic surface than the monomer and are thus retained more
strongly. HIC is also valued because it can be run directly on a high-salt
ion-exchange eluate with little conditioning, and because it elutes the product
into a low-salt buffer suited to the next step.
\end{appbox}

\begin{databox}{HIC}
The salting-out coefficient $K_{s,i}$ and prefactor $\beta_i$ are obtained from
\emph{isocratic retention measurements at several salt concentrations}: a plot of
$\ln k'$ against salt molarity is linear with slope $\propto K_{s,i}$ and
intercept $\propto \ln\beta_i$. The choice and concentration range of the salt
are themselves design variables (a Hofmeister scouting study). Saturation
capacity $q_{\max,i}$ again comes from frontal-analysis breakthrough curves at
the load (high-salt) condition.
\end{databox}

\begin{limitbox}{HIC}
The single-exponential law~\eqref{eq:hic} is an empirical, mean-field
representation: it captures the dominant log-linear salt dependence but not the
curvature seen at very high salt, nor the protein conformational changes that HIC
ligands can induce, nor any explicit dependence on temperature or salt type
(both of which are folded into $K_{s,i}$ and must be re-calibrated if they
change). As with the other modes, $q_{\max,i}$ is taken constant.
\end{limitbox}

%==============================================================================
\section{Reversed-phase chromatography (RP-HPLC)}
\label{sec:rp}

In reversed-phase chromatography the stationary phase is strongly hydrophobic
(e.g.\ alkyl-bonded silica) and the modulator is the volume fraction $\varphi$ of
an organic solvent (acetonitrile, methanol) in a predominantly aqueous mobile
phase. Solutes partition more into the mobile phase as $\varphi$ rises, so
binding falls and elution proceeds by an \emph{increasing} organic gradient. The
package uses the \emph{linear-solvent-strength} (LSS) law (class
\code{LinearSolventStrengthLaw}):
%
\begin{equation}
b_i(\varphi) = \beta_i \, \exp\!\bigl(-S_i\,\varphi\bigr),
\label{eq:lss}
\end{equation}
%
where $\beta_i$ is the affinity in purely aqueous mobile phase ($\varphi=0$) and
$S_i$ is the \emph{solvent-strength} slope. Equation~\eqref{eq:lss} is the
exponential form of the celebrated empirical observation that
$\log k' = \log k'_0 - S\varphi$ over a useful range of $\varphi$.

\begin{center}
\renewcommand{\arraystretch}{1.3}
\begin{tabular}{@{}l l p{0.55\linewidth}@{}}
\toprule
\textbf{Symbol} & \textbf{Units} & \textbf{Meaning} \\
\midrule
$\beta_i$ & --- & Affinity in pure aqueous eluent (\code{beta}). \\
$S_i$ & --- & Solvent-strength parameter; the slope of $\ln b_i$ with organic
   fraction. For large biomolecules $S_i$ is large, which is why peptides and
   proteins elute over a narrow $\varphi$ window (\code{s}). \\
$q_{\max,i}$ & g\,L$^{-1}$ & Saturation capacity (\code{q\_max}). \\
$\varphi$ & --- & Modulator: organic volume fraction, $0\le\varphi\le 1$. \\
\bottomrule
\end{tabular}
\end{center}

\begin{appbox}{reversed-phase chromatography}
RP-HPLC offers the highest resolving power of the common modes and is the
backbone of analytical characterisation---peptide mapping, purity and
impurity profiling, and the separation of closely related variants. Preparatively
it is used for peptides, oligonucleotides, and some small proteins, and for
insulin and similar molecules at manufacturing scale. Its resolving power comes
at the cost of organic solvents and, often, conditions that can denature larger
proteins---hence its dominance in analysis and small-molecule/peptide
purification rather than in the gentle capture of fragile biologics.
\end{appbox}

\begin{databox}{RP-HPLC}
The LSS parameters $S_i$ and $\beta_i$ are obtained from a handful of
\emph{linear-gradient scouting runs} at different gradient slopes (or isocratic
runs at different $\varphi$): for each solute the retention is fitted to
$\log k' = \log\beta_i + \log q_{\max,i} - S_i\varphi$. Two well-chosen gradient
runs suffice in principle to fix the pair $(S_i,\beta_i)$ for each solute---this
two-run calibration is the basis of widely used method-development software.
\end{databox}

\begin{limitbox}{RP-HPLC}
The LSS law is accurate over a limited $\varphi$ range; real $\log k'$ versus
$\varphi$ plots are gently curved (quadratic), and the linear form breaks down
near the extremes of composition. The model represents neither the temperature
dependence of retention, the influence of ion-pairing or mobile-phase additives,
nor on-column conformational effects, all of which are subsumed into effective
constants that must be re-measured if conditions change.
\end{limitbox}

%==============================================================================
\section{Multi-component resolution and the nonlinear regime}
\label{sec:multicomp}

The linear (Henry) form~\eqref{eq:henry} suffices to predict \emph{where} a
dilute band elutes, but resolution of closely related species on a
\emph{loaded} column is a nonlinear phenomenon: as the bound phase fills, the
competition denominator of~\eqref{eq:langmuir} couples the components, and a more
strongly bound species can displace a weaker one. The builder
\code{high\_resolution\_iex} assembles exactly this case---an SMA ion exchanger
run with finite $q_{\max,i}$ and full competition---so that the interplay of
selectivity (differences in $\nu_i$ and $\beta_i$), loading, and gradient slope
can be studied. The same competition is what produces the self-sharpening fronts
and the displacement effects of overloaded preparative chromatography.

The single most important operating variable for resolution is the
\emph{gradient slope}. Figure~\ref{fig:slope} demonstrates the universal
trade-off it controls: a steep gradient elutes the band early, tall, and narrow,
but compresses the separation in modulator space (poor selectivity between
neighbours); a shallow gradient spreads the separation out, improving resolution
at the cost of a later, broader, more dilute peak and a longer cycle.

\begin{figure}[t]
\centering
\includegraphics[width=0.9\linewidth]{gradient_slope.png}
\caption[Gradient slope and the resolution--dilution trade-off]{The same protein
eluted under three salt-gradient slopes (dotted: the gradients), computed by
\code{run\_column}; $N$ is the apparent plate count from
Eq.~\eqref{eq:plates}. A steep gradient (green) gives an early, sharp,
concentrated peak; a shallow gradient (red) gives a late, broad, dilute one. This
single knob is the principal lever of preparative method development, and a
natural factor for the experimental designs of Chapter~\ref{ch:design}.}
\label{fig:slope}
\end{figure}

%==============================================================================
\section{The elution program: specifying a gradient}
\label{sec:program}

A run is described (\code{models/chromatography/program.py}) as an ordered list
of \emph{segments}, each measured in \emph{column volumes} (CV) and each either
holding the modulator constant or ramping it linearly. A segment is a
\code{Segment(name, duration\_cv, m\_start, m\_end)}; a linear ramp---the
``linearly changing eluate'' that is the gradient---is the case
$m_{\mathrm{end}}\neq m_{\mathrm{start}}$. A typical bind-and-elute program is
the five-segment sequence built by \code{ElutionProgram.linear\_gradient}:
%
\[
\text{equilibrate} \;\to\; \text{load} \;\to\; \text{wash}
\;\to\; \text{gradient} \;\to\; \text{strip}.
\]

Durations in CV are converted to seconds against a particular column and flow.
One column volume passes in
%
\begin{equation}
t_{\mathrm{CV}} = \frac{V_{\mathrm{col}}}{Q} = \frac{V_{\mathrm{col}}}{u A \eps}
= \frac{L}{u\,\eps},
\label{eq:tcv}
\end{equation}
%
after which the inlet modulator $m_{\mathrm{in}}(t)$ is the piecewise-linear
function interpolated across the segment breakpoints. The protein
\code{Injection} is a separate rectangular feed pulse. Its duration follows from
the \emph{load density} $\rho_{\mathrm{load}}$ (grams of protein per litre of
adsorbent): the volume that delivers it is $V_{\mathrm{inject}} =
\rho_{\mathrm{load}} V_{\mathrm{resin}} / c_{\mathrm{feed}}$, which expressed in
column volumes is
%
\begin{equation}
\mathrm{CV}_{\mathrm{load}}
= \frac{\rho_{\mathrm{load}}\,(1-\eps)}{c_{\mathrm{feed}}},
\label{eq:loadcv}
\end{equation}
%
with $c_{\mathrm{feed}}$ the total feed concentration
(\code{Injection.from\_load\_density}). Load density, gradient slope, load and
elution pH, and start/end salt are thus the natural \emph{critical process
parameters} that the designs of Chapter~\ref{ch:design} vary.

%==============================================================================
\section{Performance and separation metrics}
\label{sec:metrics}

A simulated chromatogram is only useful once reduced to the quantities a process
is specified on. The package (\code{models/chromatography/metrics.py}) computes
two families of metric.

\paragraph{Pool metrics.}
Given a collection window (the ``cut points'') $[t_1,t_2]$, the outlet
chromatogram of the target component is integrated by the trapezoidal rule to
give the pooled mass, and three ratios follow:
%
\begin{align}
\text{yield} &= \frac{\int_{t_1}^{t_2} c_{\mathrm{tgt}}\,\dd t}
                      {\int_{0}^{\infty} c_{\mathrm{tgt}}\,\dd t}, \\[0.3em]
\text{purity} &= \frac{\int_{t_1}^{t_2} c_{\mathrm{tgt}}\,\dd t}
                      {\sum_i \int_{t_1}^{t_2} c_i\,\dd t}, \\[0.3em]
\text{productivity} &= \frac{\int_{t_1}^{t_2} c_{\mathrm{tgt}}\,\dd t}
                            {t_{\mathrm{cycle}}}.
\end{align}
%
Yield, purity, and productivity are the canonical, and usually competing,
objectives of a downstream step---and exactly the responses optimised by the
Bayesian loop of Chapter~\ref{ch:design}.

\paragraph{Peak metrics.}
For peak shape the package uses \emph{statistical moments}, which are robust to
the skew of gradient peaks where a Gaussian fit would be biased. With the
zeroth, first, and second moments giving the area, the retention time, and the
variance,
%
\begin{equation}
A = \!\int\! c\,\dd t, \quad
t_R = \frac{1}{A}\!\int\! t\,c\,\dd t, \quad
\sigma^2 = \frac{1}{A}\!\int\! (t-t_R)^2 c\,\dd t,
\end{equation}
%
the apparent \emph{plate count} and the \emph{resolution} between two peaks are
%
\begin{equation}
N = \frac{t_R^2}{\sigma^2}, \qquad
R_s = \frac{2\,|t_{R,b}-t_{R,a}|}{w_a + w_b}, \quad w = 4\sigma.
\label{eq:plates}
\end{equation}
%
The plate count $N$ measures efficiency (how narrow a band the column produces);
the resolution $R_s$ measures how cleanly two bands are separated, with
$R_s\gtrsim 1.5$ the usual target for baseline separation.

%==============================================================================
\section{The reduced equilibrium-dispersive model}
\label{sec:legacy}

For fast single-component, isocratic work the package retains the original
equilibrium-dispersive model (\code{models/chromatography/legacy.py}). Here the
fast-equilibrium closure is used: $q = H(m,\mathrm{pH})\,c$ is substituted into
the mass balance, and the time derivative of the bound phase is folded into a
\emph{retardation factor}
%
\begin{equation}
R(m,\mathrm{pH}) = 1 + \frac{1-\eps}{\eps}\,H(m,\mathrm{pH}),
\label{eq:retard}
\end{equation}
%
leaving a single scalar transport equation per component,
%
\begin{equation}
R\,\pdv{c}{t} = -\,u\,\pdv{c}{z} + D_{\mathrm{ap}}\,\pdv[2]{c}{z}.
\label{eq:ed}
\end{equation}
%
The retardation factor is the ratio of the velocity of the unretained solvent to
that of the protein band: a band with $R=100$ moves through the column a hundred
times more slowly than the mobile phase. Because $H$ for a strongly retained
protein can reach $10^4$, $R$ can be enormous, which is the origin of the
stiffness noted in Section~\ref{sec:engine} and the reason the integrator is
again BDF. The ED model uses the same linearised SMA Henry
constant~\eqref{eq:smahenry} and is restricted to a constant modulator; it is
faster than the full engine precisely because it carries one field per component
instead of two plus the modulator.

%==============================================================================
\section{Resolving the mass-transfer mechanism: the general rate model}
\label{sec:grm}

The transport-dispersive engine of Section~\ref{sec:engine} lumps every
mass-transfer resistance into a single linear-driving-force coefficient
$k_{m,i}$. That is adequate when one only wants band positions and an effective
plate count, but it conflates two physically distinct resistances---the external
\emph{film} (boundary layer) around each bead and the \emph{pore diffusion}
inside it---and, being a stiff lumped sink, it can defeat the method-of-lines
integrator: for strongly bound bands under steep gradients the adaptive
\textsc{bdf} solver can take a step that drives the protein band negative and
then to zero, \emph{silently losing the entire injected mass}
(Figure~\ref{fig:grm-cons}, red). The package therefore provides a second,
independent solver---a \emph{general rate model} (GRM) discretised by the finite
volume method through the \code{PyFVTool} library
(\code{models/chromatography/grm.py})---that both resolves the mechanism and is
mass-conservative by construction.

\subsection{Governing equations}

The GRM keeps the bulk (interstitial) balance of Section~\ref{sec:engine} but now
couples it, through a film boundary layer, to a \emph{resolved} radial diffusion
problem inside a representative spherical bead at each axial location. For each
component $i$,
%
\begin{align}
\text{bulk:}\quad
&\pdv{c_i}{t} = -u\,\pdv{c_i}{z} + D_{\mathrm{ax}}\,\pdv[2]{c_i}{z}
  - \frac{1-\eps}{\eps}\,\frac{3}{R_p}\,k_{f,i}\,
    \bigl(c_i - c_{p,i}\big|_{r=R_p}\bigr),
\label{eq:grm-bulk}\\[0.3em]
\text{particle:}\quad
&\eps_p\,\pdv{c_{p,i}}{t} + (1-\eps_p)\,\pdv{q_i}{t}
  = \eps_p D_{p,i}\,\frac{1}{r^2}\pdv{}{r}\!\left(r^2 \pdv{c_{p,i}}{r}\right),
\label{eq:grm-part}
\end{align}
%
with the adsorbed loading $q_i = \qstar_i(\bm c_p, m, \mathrm{pH})$ given by the
\emph{same} isotherm~\eqref{eq:langmuir} as the lumped engine, so every mode and
the competitive multi-component case carry over unchanged. The two new boundary
conditions on the bead express film continuity at the surface and symmetry at the
centre,
%
\begin{equation}
\eps_p D_{p,i}\left.\pdv{c_{p,i}}{r}\right|_{r=R_p}
  = k_{f,i}\bigl(c_i - c_{p,i}\big|_{r=R_p}\bigr),
\qquad
\left.\pdv{c_{p,i}}{r}\right|_{r=0} = 0.
\label{eq:grm-bc}
\end{equation}
%
The modulator $m$ is an unretained tracer; it is advanced on the bulk mesh and
taken to equilibrate instantly within the pores (pore modulator $=$ bulk
modulator at that axial location), which decouples it from the protein solve.

\subsection{Mechanistic mass-transfer parameters}

The two resistances are given physical, correlation-based defaults
(\code{film\_coefficient}, overridable by direct input). The \emph{film}
coefficient follows the Wilson--Geankoplis packed-bed correlation,
%
\begin{equation}
\eps\,\mathrm{Sh} = 1.09\,\mathrm{Re}^{1/3}\mathrm{Sc}^{1/3},
\qquad
k_{f} = \frac{\mathrm{Sh}\,D_m}{d_p},
\label{eq:wilson}
\end{equation}
%
with $\mathrm{Re}=\rho u_s d_p/\mu$ (superficial velocity $u_s=\eps u$),
$\mathrm{Sc}=\mu/\rho D_m$, $d_p=2R_p$, and $D_m$ the free-solution diffusivity.
The external specific area is $a_v = 3(1-\eps)/R_p$. The intraparticle transport
is set by the \emph{pore diffusivity} $D_{p,i}$ (defaulting to
$\eps_p D_m/\tau$ with tortuosity $\tau$). Together these make
``boundary-layer'' and ``inside-the-bead'' resistances separately tunable and
physically grounded, and Figure~\ref{fig:grm-prof} shows the radial profiles the
model resolves---a steep surface-to-centre gradient that the lumped model cannot
represent.

\begin{figure}[t]
\centering
\includegraphics[width=0.78\linewidth]{grm_intraparticle.png}
\caption[Resolved intraparticle profiles]{Pore-phase concentration $c_p(r)$
inside the bead at four column positions, during loading, computed by
\code{run\_grm}. Each bead shows the surface-to-centre gradient set by the film
and pore-diffusion resistances; beads deeper in the column (toward the outlet)
are less loaded. This radial structure is the mechanistic content the general
rate model adds over the lumped LDF engine.}
\label{fig:grm-prof}
\end{figure}

\subsection{Finite-volume discretisation and exact conservation}

The bulk axial operators (upwind convection, central dispersion) are built on a
Cartesian grid and the radial operator on a spherical grid, both with
\code{PyFVTool}; the bulk and the per-axial-cell bead problems are assembled into
one global sparse system and advanced \emph{fully implicitly} (backward Euler),
with the competitive isotherm linearised by Picard iteration each step. Two
details secure machine-precision conservation. First, the bulk$\leftrightarrow$
bead film exchange is written so the \emph{identical} discrete flux appears as a
sink in the bulk cell and a source in the bead surface cell. Second---a subtle
point---\code{PyFVTool}'s spherical \code{cellvolume} is the midpoint
approximation $4\pi r_c^2\Delta r$, whereas its diffusion operator conserves
against the \emph{true} shell volume $\tfrac43\pi(r_{\mathrm{out}}^3 -
r_{\mathrm{in}}^3)$; the GRM uses the true shell volume for all mass accounting
and film normalisation. Finally, the adsorption is discretised in
\emph{storage form},
%
\begin{equation}
\frac{\bigl[\eps_p c_{p,i} + (1-\eps_p)\qstar_i(\bm c_p,m)\bigr]^{n+1}
      - \bigl[\eps_p c_{p,i} + (1-\eps_p)\qstar_i(\bm c_p,m)\bigr]^{n}}{\Delta t}
= \eps_p D_{p,i}\nabla_r^2 c_{p,i}^{\,n+1},
\label{eq:grm-storage}
\end{equation}
%
i.e.\ the change in \emph{stored} mass is differenced exactly rather than through
a Jacobian increment. Because the old storage is evaluated at the old modulator,
this single device captures both the nonlinear isotherm \emph{and} the changing
gradient with no separate $\dd q/\dd t$ source term, and conserves mass to
machine precision.

\subsection{Verification and comparison}

On a benign gradient where the lumped engine is reliable, the two solvers agree
closely (Figure~\ref{fig:grm-val}), validating the GRM. On a stiff,
strongly-bound gradient the difference is decisive (Figure~\ref{fig:grm-cons}):
the method-of-lines engine loses $100\%$ of the injected mass (its outlet is
identically zero), while the GRM conserves to a relative closure error of
$\sim\!10^{-11}$ and returns the physically correct eluting band. This is the
regime---competitive, strongly-retained, multi-component gradient
elution---that motivated the second solver.

\begin{figure}[t]
\centering
\includegraphics[width=0.82\linewidth]{grm_vs_mol_validation.png}
\caption[GRM/MoL validation]{On a benign single-component gradient the
method-of-lines (lumped LDF) and general rate model solvers produce essentially
the same chromatogram, validating the finite-volume GRM against the established
engine.}
\label{fig:grm-val}
\end{figure}

\begin{figure}[t]
\centering
\includegraphics[width=\linewidth]{grm_conservation.png}
\caption[GRM/MoL mass conservation]{A stiff, strongly-bound gradient.
\emph{Left:} the method-of-lines engine's band vanishes---it loses all the
injected mass---while the GRM returns a proper eluting peak. \emph{Right:} the
mass-balance closure error (log scale): $\sim\!100\%$ for the MoL engine in this
regime versus $\sim\!10^{-11}$ for the conservative finite-volume GRM.}
\label{fig:grm-cons}
\end{figure}

\begin{remark}
The GRM is the more faithful and more robust model, at a higher cost: it carries
$n_c(N_z + N_z N_r)$ unknowns plus the modulator and solves a coupled implicit
system each step. The lumped engine remains the right tool for fast design-space
sweeps in benign regimes; the GRM is the reference for mechanism studies, for
mass-transfer-limited separations, and wherever the lumped solver's conservation
is in doubt.
\end{remark}

%==============================================================================
\section{Consolidated limits of validity}
\label{sec:limits-summary}

\begin{limitbox}{the model family as a whole}
\begin{itemize}
  \item \textbf{Lumped dispersion and mass transfer.} Pore diffusion, film
  resistance, and eddy dispersion are not resolved individually; they are
  represented by the apparent dispersion $D_{\mathrm{ap}}$ and the LDF
  coefficient $k_{m,i}$. The model reproduces band \emph{shapes} and
  \emph{positions}, not the microscopic transport that underlies them.
  \item \textbf{Isothermal, single-phase, one-dimensional.} The column is taken
  as radially uniform and at constant temperature; wall effects, channelling,
  and thermal gradients are absent.
  \item \textbf{Constitutive simplifications.} Each affinity law is the dominant
  empirical dependence of its mode, calibrated over a finite window; saturation
  capacities are constant, and steric, conformational, and Hofmeister effects
  are absorbed into effective constants.
  \item \textbf{Numerical caveats.} First-order upwinding adds grid dispersion,
  so the cell count participates in setting band width; and for very strong
  binding under steep gradients the stiff integration can lose mass for the
  most-retained bands unless dispersion and solver tolerances are chosen with
  care. The robust regime is moderate binding with adequate dispersion.
\end{itemize}
Within these bounds the models are fit for their purpose: to serve as a fast,
differentiable, physically faithful oracle whose response to the critical
process parameters has the right shape and the right sensitivities---which is all
that the experimental-design methods of the next chapter require.
\end{limitbox}
